
% CREACIÓN DEL DOCUMENTO
\documentclass[
	spanish, % Idioma: spanish, english, etc.
	letterpaper, oneside
]{article}

% INFORMACIÓN DEL DOCUMENTO
\def\documenttitle {Cuadro comparativo de escuelas éticas}
\def\documentsubtitle {Escuelas éticas}
\def\documentsubject {Actividad 1 - Ética y Moral jurídica}

\def\documentauthor {Martín Jonathan de la Cruz}
\def\coursename {Ética y Moral jurídica}
\def\coursecode {030}

\def\universityname {Universidad Abierta y a Distancia de México}
\def\universityfaculty {Licenciatura en Derecho}
\def\universitydepartment {Ética y Moral jurídica}
\def\universitydepartmentimage {departamentos/UnADM}
\def\universitydepartmentimagecfg {height=1.57cm}
\def\universitylocation {Roma Norte, Ciudad de México}

% INTEGRANTES, PROFESORES Y FECHAS
\def\authortable {
	\begin{tabular}{ll}
		\textbf{Alumno:}
		& \begin{tabular}[t]{l}
			 Martín Jonathan de la Cruz \\	
		\end{tabular} \\
		Matrícula: &   ES2611202040 \\
		
		\textit{Profesor:}
		& \begin{tabular}[t]{l}
			Emma Liliana \\ Padilla Cano
		\end{tabular} \\
		& \\
		\multicolumn{2}{l}{Fecha de realización: \today} \\
		% \multicolumn{2}{l}{Fecha de entrega: \today} \\
		\multicolumn{2}{l}{\universitylocation}
	\end{tabular}
}

% IMPORTACIÓN DEL TEMPLATE
\input{template}

% FORMATO SOLICITADO POR LA PLANTILLA
\usepackage{helvet}
\renewcommand{\familydefault}{\sfdefault}

% MARCA DE AGUA TIKZ EN PORTADA
\usepackage{tikz}

% CITAS EN FORMATO AUTOR-AÑO (APA)
\setcitestyle{authoryear,open={(},close={)}}

% ==============================================================================
% MARCA DE AGUA EN PORTADA
% - Se aplica solo a la portada porque usa \AddToShipoutPictureBG*.
% - Para apagarla en alguna actividad: \def\coverwatermarkenabled {false}
% - Usar opacidad baja para que no compita con titulo, datos ni logotipo.
% ==============================================================================
\def\coverwatermarkenabled {true}
\def\coverwatermarkimage {img/departamentos/UnADM.pdf}
\def\coverwatermarkopacity {0.16}
\def\coverwatermarkwidth {0.72\paperwidth}
\def\coverwatermarkxshift {0cm}
\def\coverwatermarkyshift {-0.35cm}

\newcommand{\insertcoverwatermark}{%
	\ifthenelse{\equal{\coverwatermarkenabled}{true}}{%
		\AddToShipoutPictureBG*{%
			\begin{tikzpicture}[remember picture,overlay]
				\node[
					opacity=\coverwatermarkopacity,
					inner sep=0pt,
					xshift=\coverwatermarkxshift,
					yshift=\coverwatermarkyshift
				] at (current page.center) {%
					\includegraphics[width=\coverwatermarkwidth]{\coverwatermarkimage}%
				};
			\end{tikzpicture}%
		}%
	}{}%
}

% INICIO DE PÁGINAS
\begin{document}

% PORTADA
\insertcoverwatermark
\templatePortrait

% CONFIGURACIÓN DE PÁGINA Y ENCABEZADOS
\templatePagecfg
\onehalfspacing

% RESUMEN O ABSTRACT
\begin{abstractd}

	\textbf{Objetivo:} Analizar y comparar las principales escuelas clásicas de la ética mediante un cuadro comparativo, con el fin de identificar sus fundamentos sobre el bien, la virtud, el deber y su relevancia para la formación del hombre jurídico.

	\textbf{Producto elaborado:} Se presenta un cuadro comparativo sobre la ética de la virtud, la deontología, el hedonismo, el estoicismo, el epicureísmo y el cinismo, organizado a partir de categorías comunes que permiten distinguir representantes, fundamentos morales, fines humanos, criterios de acción y aportes para el ámbito jurídico.

	\textbf{Enfoque de análisis:} El trabajo parte de la idea de que cada escuela ética propone una forma distinta de comprender la vida buena y de orientar la conducta humana. Por ello, la comparación no solo describe conceptos, sino que también permite valorar cómo estas corrientes pueden aportar criterios útiles para la formación ética del profesional del derecho \citep{huerta2000etica,ronquillo2018etica,prieto2009favor}.

\end{abstractd}

% TABLA DE CONTENIDOS - ÍNDICE
\templateIndex

% CONFIGURACIONES FINALES
\templateFinalcfg

% ======================= INICIO DEL DOCUMENTO =======================

\section{Introducción}

La ética, como rama fundamental de la filosofía, se encarga del estudio de la conducta humana y de los principios que orientan la distinción entre lo correcto y lo incorrecto \citep{ronquillo2018etica}. A lo largo de la historia, diversas escuelas filosóficas han propuesto modelos explicativos sobre cómo debe vivir el ser humano y cuáles son los fines que deben guiar sus acciones.

En este contexto, el análisis comparativo de las principales corrientes éticas permite comprender no solo sus postulados individuales, sino también las tensiones y coincidencias entre ellas. La ética de la virtud, la deontología, el hedonismo, el estoicismo, el epicureísmo y el cinismo representan distintas formas de concebir la moralidad, el bienestar y la libertad \citep{huerta2000etica}.

El presente trabajo tiene como propósito desarrollar un cuadro comparativo que facilite la organización, análisis y contrastación de estas corrientes, promoviendo una visión integral y crítica del pensamiento ético clásico.

\section{Escuelas éticas clásicas}

Se comparan las siguientes escuelas éticas clásicas: ética de la virtud (Aristóteles), deontológica (Kant), hedonismo, estoicismo (Séneca y Epicteto), epicureísmo (Epicuro) y cinismo (Diógenes).

\subsection*{Ética de la virtud (Aristóteles)}

La ética de la virtud sostiene que la felicidad (\textit{eudaimonía}) es el fin último del ser humano y se alcanza mediante el desarrollo de hábitos virtuosos guiados por la razón. La virtud consiste en el justo medio entre extremos, evitando tanto el exceso como el defecto. La vida buena es una actividad racional conforme a la virtud \citep{huerta2000etica}.

\subsection*{Ética deontológica (Kant)}

La ética deontológica establece que la moralidad depende del cumplimiento del deber conforme a la ley moral. Según Kant, una acción es moral cuando se realiza por respeto a la ley y no por inclinaciones personales. El valor moral reside en la buena voluntad y en la intención del sujeto \citep{huerta2000etica}.

\subsection*{Hedonismo}

El hedonismo sostiene que el fin último de la vida es el placer y la ausencia de dolor. Considera que lo bueno es aquello que genera satisfacción y lo malo lo que produce sufrimiento. Esta corriente coloca el placer sensible como valor central de la conducta humana \citep{huerta2000etica}.

\subsection*{Estoicismo}

El estoicismo afirma que la virtud es el único bien verdadero y que la felicidad depende del dominio interior. El ser humano debe enfocarse en aquello que puede controlar y aceptar con serenidad lo que no depende de él. La vida virtuosa implica autodisciplina y racionalidad \citep{huerta2000etica}.

\subsection*{Epicureísmo}

El epicureísmo plantea que el placer es el bien supremo, pero entendido como un placer moderado y racional. Su objetivo es alcanzar la tranquilidad del alma (\textit{ataraxia}), evitando excesos y sufrimientos innecesarios. Promueve una vida sencilla y equilibrada \citep{huerta2000etica}.

\subsection*{Cinismo}

El cinismo propone una vida conforme a la naturaleza, basada en la autosuficiencia y el rechazo de las convenciones sociales. Defiende una existencia austera y libre de dependencias materiales o sociales, privilegiando la autenticidad y la independencia personal \citep{huerta2000etica}.

\section{Tabla comparativa de escuelas éticas clásicas}

% \subsection*{Proceso de elaboración del cuadro comparativo}

% \begin{enumerate}
% 	\item \textbf{Seleccionar.} Subrayar, separar y catalogar la información relevante y útil de cada escuela ética para contar con insumos suficientes para la comparación.
% 	\item \textbf{Trazar la tabla.} Definir el número de conceptos o temas y la cantidad de categorías que se usarán en la comparación. Cada concepto se coloca en una columna y cada categoría en una fila.
% 	\item \textbf{Titular el cuadro.} Enunciar con claridad los conceptos o temas que se comparan.
% 	\item \textbf{Distribuir la información.} Colocar los conceptos en la parte superior de las columnas y ubicar las categorías en las celdas correspondientes, con nombres puntuales y claros.
% 	\item \textbf{Revisar.} Verificar la pertinencia de la información, la redacción y la claridad general del cuadro comparativo.
% \end{enumerate}

\clearpage
\pagestyle{empty}
\begin{landscape}
\thispagestyle{empty}
\begingroup
\scriptsize
\setlength{\tabcolsep}{4pt}
\renewcommand{\arraystretch}{1.35}
\begin{longtable}{@{}>{\raggedright\arraybackslash}p{0.16\linewidth}>{\raggedright\arraybackslash}p{0.13\linewidth}>{\raggedright\arraybackslash}p{0.13\linewidth}>{\raggedright\arraybackslash}p{0.13\linewidth}>{\raggedright\arraybackslash}p{0.13\linewidth}>{\raggedright\arraybackslash}p{0.13\linewidth}>{\raggedright\arraybackslash}p{0.13\linewidth}@{}}
\caption{Cuadro comparativo de escuelas éticas clásicas}\label{tab:cuadro-eticas}\\
\toprule
\textbf{Categoría} & \textbf{Ética de la virtud} & \textbf{Deontológica} & \textbf{Hedonismo} & \textbf{Estoicismo} & \textbf{Epicureísmo} & \textbf{Cinismo} \\
\midrule
\endfirsthead
\toprule
\textbf{Categoría} & \textbf{Ética de la virtud} & \textbf{Deontológica} & \textbf{Hedonismo} & \textbf{Estoicismo} & \textbf{Epicureísmo} & \textbf{Cinismo} \\
\midrule
\endhead
\bottomrule
\endfoot
\bottomrule
\endlastfoot
Representante principal & Aristóteles & Immanuel Kant & Aristipo de Cirene (tradición hedonista) & Séneca y Epicteto & Epicuro & Diógenes de Sinope \\
\midrule
Fundamento moral & Formación del carácter por hábitos virtuosos & Cumplimiento del deber y ley moral universal & Búsqueda del placer y evitación del dolor & Virtud como único bien verdadero & Placer racional y moderado & Vida natural y autosuficiente \\
\midrule
Fin último del ser humano & Eudaimonía (vida lograda) & Obrar por deber con buena voluntad & Obtener placer y minimizar sufrimiento & Alcanzar serenidad interior mediante dominio de sí & Ataraxia (tranquilidad del alma) & Libertad frente a convenciones sociales \\
\midrule
Criterio para actuar bien & Justo medio guiado por la razón & Imperativo categórico & Consecuencias placenteras de la acción & Conformidad con la razón y la naturaleza & Evaluación prudente de placeres y dolores & Rechazo de lo superfluo y lo artificial \\
\midrule
Actitud ante bienes externos & Son útiles, pero subordinados a la virtud & No determinan el valor moral & Son valiosos si producen placer & Son indiferentes frente a la virtud & Se aceptan con moderación & Se minimizan para preservar la independencia \\
\midrule
Aporte para el hombre jurídico & Énfasis en prudencia, justicia y formación ética profesional & Centralidad del deber, norma y universalidad & Consideración del bienestar y daño en decisiones & Autocontrol, resiliencia y rectitud ante la adversidad & Moderación y equilibrio en la toma de decisiones & Crítica de prácticas jurídicas vacías y apego a la autenticidad \\
\end{longtable}
\endgroup
\end{landscape}
\clearpage
\pagestyle{fancy}

La comparación evidencia que todas las escuelas buscan orientar la acción humana hacia una vida buena, pero difieren en el criterio rector: virtud, deber, placer, autodominio, moderación o autosuficiencia. Esta distinción permite valorar sus aportes para la formación ética del profesional del derecho \citep{huerta2000etica,prieto2009favor}.


\section{Conclusión}

El análisis comparativo de las escuelas éticas clásicas pone de relieve la diversidad de enfoques sobre la moralidad y el bien. Cada corriente ofrece una perspectiva particular sobre cómo alcanzar la felicidad y vivir de manera virtuosa, lo que enriquece la comprensión de la ética como disciplina. Sin embargo, todas coinciden en la necesidad de orientar la conducta humana hacia un ideal de bien y realización personal. Desde la formación jurídica, esta revisión resulta valiosa porque permite reconocer que la ética no solo acompaña al Derecho como reflexión teórica, sino que también ofrece criterios para interpretar normas, asumir responsabilidades y actuar con mayor sentido de justicia en la práctica profesional \citep{ronquillo2018etica,huerta2000etica}.\footnote{En la elaboración de este trabajo se utilizaron herramientas de inteligencia artificial como apoyo para la organización y revisión del texto, conforme a los Lineamientos para el uso ético y responsable de la Inteligencia Artificial en el ámbito académico de la UnADM. El análisis, la selección de fuentes y la postura sostenida son de mi autoría.}

\bibliography{etica-y-moral-juridica}

% FIN DEL DOCUMENTO
\end{document}

